% Generated draft: requires mathematical and source review.
\documentclass[11pt]{article}
\usepackage[margin=25mm]{geometry}
\usepackage{fontspec}
\setmainfont{Latin Modern Roman}
\usepackage{amsmath,amssymb,mathtools,unicode-math}
\setmathfont{Latin Modern Math}
\usepackage{enumitem,needspace,xurl,hyperref}
\setlength{\parindent}{0pt}
\setlength{\parskip}{5pt}
\setlength{\emergencystretch}{3em}
\newcommand{\N}{\mathbb{N}}
\newcommand{\Z}{\mathbb{Z}}
\newcommand{\Q}{\mathbb{Q}}
\newcommand{\R}{\mathbb{R}}
\newcommand{\C}{\mathbb{C}}
\newcommand{\F}{\mathbb{F}}
\newcommand{\A}{\mathbb{A}}
\newcommand{\PP}{\mathbb{P}}
\newcommand{\E}{\mathbb{E}}
\newcommand{\eps}{\varepsilon}
\newcommand{\dd}{\,\mathrm{d}}
\begin{document}
\section*{UnsolvedMath ID: 30003354}
\textbf{Source problem code:} catalogue entry
\subsection{Definitions and mathematical statement}
Let \(V\) be one of the simply connected surfaces \(\R^2\) and \(S^2\). The supplied record concerns conformal metric parametrizations of the form
\[
F_V(u,\phi)=\phi^*\bigl(e^{-2u}g_\kappa\bigr),
\]
where \(g_\kappa\) denotes the reference Riemannian metric appearing in the source, \(u\) is a real-valued function on \(V\), and \(\phi\) is a diffeomorphism of \(V\), with sufficient regularity for the expression to be defined. Explicitly, for \(x\in V\) and \(v,w\in T_xV\),
\[
F_V(u,\phi)_x(v,w)
=e^{-2u(\phi(x))}(g_\kappa)_{\phi(x)}
\bigl(\dd\phi_x(v),\dd\phi_x(w)\bigr).
\]
Write \(\mathcal U_V\), \(\mathcal D_V\), and \(\mathcal M_V\) as placeholders for the source's intended admissible function space, diffeomorphism space, and metric space, respectively; these spaces are not specified in the supplied record.

The requested borderline question is whether, for each \(V\in\{\R^2,S^2\}\), the intended parametrization
\[
F_V:\mathcal U_V\times\mathcal D_V\longrightarrow\mathcal M_V
\]
is a homeomorphism when the function, diffeomorphism, and metric spaces carry the intended \(C^{r,0}\) topologies, and the domain carries the corresponding product topology. A homeomorphism here means a bijective continuous map whose inverse is continuous.

\textbf{Unresolved formulation qualifications.} The record does not define \(g_\kappa\) or \(\kappa\), the admissible spaces, the range of \(r\), or the precise topologies and differentiability orders on those spaces. It does not state the normalization of the conformal factors and diffeomorphisms needed to specify the claimed parametrization. Although the supplied references concern nonnegatively curved surfaces, the record does not give the exact curvature, completeness, or other admissibility conditions for \(\mathcal M_V\), nor the corresponding restrictions on \(\mathcal U_V\). On the noncompact surface \(\R^2\), it also leaves unspecified any conditions at infinity and whether the topology uses convergence on compact subsets or global control. Thus the displayed question preserves the supplied target but remains schematic pending recovery of the missing source hypotheses; no particular choices of these data are asserted.

The supplied dated curation reports that the decomposition map is a homeomorphism in the source's stated \(C^{r,\alpha}\) settings for \(0<\alpha<1\), while the borderline \(\alpha=0\) case remains unresolved and is expected to fail. This is a supplied literature assessment, not an independently verified theorem or status determination.
\subsection{Short English statement}
For the plane and the sphere, ask whether the source's conformal parametrization of metrics remains a homeomorphism at the borderline \(C^{r,0}\) regularity. The exact admissible spaces, normalizations, reference metric, and topologies must be recovered before this becomes a fully specified mathematical question. The supplied curation labels the problem partially solved and describes failure at the borderline as an expectation.
\subsection{Sources}
\begin{itemize}
\item Supplied local UnsolvedMath registry record, permanent ID \(30003354\), exact source problem code \texttt{catalogue entry}, title \emph{Borderline Continuity of Conformal Metric Parametrizations}, in the Oberwolfach Reports Open Problems collection. No individual proposer is named; the supplied proposed year is 2017.
\item Supplied report citation: \emph{Mini-Workshop: Spaces and Moduli Spaces of Riemannian Metrics (2018)}, Oberwolfach Reports, pp.\ 133--166, DOI \texttt{10.4171/owr/2017/3}. Source URL and literature-source link labeled ``Oberwolfach report on nonnegative surfaces'': \url{https://doi.org/10.4171/owr/2017/3}. The citation's parenthetical year 2018 differs from the supplied proposed year and year tag, both 2017; this discrepancy is unresolved.
\item Supplied literature-source link: ``Banakh--Belegradek spaces of nonnegatively curved surfaces,'' \url{https://arxiv.org/abs/1510.07269}.
\item Literature assessment checked at \texttt{2026-08-21}: the supplied curation labels the problem \texttt{partially\_solved}, reports a homeomorphism result in the stated \(C^{r,\alpha}\) settings for \(0<\alpha<1\), and describes the \(\alpha=0\) case as unresolved and expected to fail. This is supplied dated curation derived from the OWR record, not independently verified; its status and cited claims require independent checking.
\item The registry assigns the statement status \texttt{corrected\_verified} and reports that the official passage was checked and rewritten while preserving its mathematical content. That verification claim is supplied metadata; the references and original passage have not been independently consulted for this draft.
\end{itemize}
\end{document}
